% ============================================================
% 第五章 模型的建立与求解
% ============================================================

\section{模型的建立与求解}

\subsection{问题一的模型建立与求解}

\subsubsection{模型建立}

数学建模必然涉及不少数学公式的使用。如果希望某个公式带编号，并且在后文中引用可以参考下面的写法：
\begin{equation}
	E=mc^2
	\label{eq:energy}
\end{equation}
式\cref{eq:energy}是质能方程。

多行公式有时候希望能够在特定的位置对齐，以下是其中一种处理方法。
\begin{align}
	P & = UI \\
	& = I^2R
\end{align}

矩阵的输入也不难。
\[
\mathbf{X} = \left(
    \begin{array}{cccc}
    x_{11} & x_{12} & \ldots & x_{1n}\\
    x_{21} & x_{22} & \ldots & x_{2n}\\
    \vdots & \vdots & \ddots & \vdots\\
    x_{n1} & x_{n2} & \ldots & x_{nn}\\
    \end{array} \right)
\]

分段函数这些可以用 \verb|cases| 环境，但是它要放在数学环境里面。
\[
f(x) =
    \begin{cases}
        0 &  x \text{为无理数} ,\\
        1 &  x \text{为有理数} .
    \end{cases}
\]

\begin{theorem}
	这是一个定理。
	\label{thm:example}
\end{theorem}

\begin{proof}
	这是一个证明。
\end{proof}

\subsubsection{模型求解}

描述求解算法与步骤，可配合流程图或伪代码进行说明。

\subsection{问题二的模型建立与求解}

\subsection{问题三的模型建立与求解}
